\documentclass{article}
\usepackage{graphicx}
\usepackage{amsmath}
\usepackage{amssymb}
\usepackage{amsfonts}
\usepackage{hyperref}
\usepackage{exercises}
\usepackage{indentfirst}
\usepackage{circledsteps}
\usepackage{amsthm}
\usepackage{csquotes}
\usepackage{yhmath}
\usepackage{comment}
\usepackage{placeins}
\usepackage{geometry}
\usepackage{cleveref}
\usepackage{stmaryrd}
\usepackage{hyperref}
\usepackage{tikz}
\usepackage{tabto}
\newcommand*{\oldneg}{\mathord{\sim}}
\newcommand{\Z}{\mathbb{Z}}
\newcommand{\N}{\mathbb{N}}
\let\oldemptyset\emptyset
\newcommand{\lbc}{\llbracket}
\newcommand{\rbc}{\rrbracket}
\newcommand{\ltop}[1]{\stackrel{\leftarrow}{#1}}

\theoremstyle{definition}
\newtheorem{definition}{Definition}
\newtheorem{theorem}{Theorem}[section]
\newtheorem{corollary}{Corollary}[theorem]
\newtheorem{lemma}[theorem]{Lemma}

\setlength{\parindent}{.5in}
\geometry{margin=0.5in}
\setlength{\parskip}{12pt}
\title{Chapter 18 Notes}
\begin{document}
\date{}
\maketitle

\section*{Section 2}

\subsection*{Lemma 2.2}
Here is why $\cup_{q_0\in\mathcal{G}}S(q_0)$ is an $M$-set (and hence so is $S$). Let 
$G(x,q)$ stand for the formula
\[ (\forall y\in x)(\textrm{Ordinal}(y)\wedge(\exists g\in R^\mathcal{G}_\beta)(q\vDash\lbc\hat{y}\approx g\rbc)) \]
(this can be expanded to a real first-order formula by results in the previous chapter).
If $S(q_0)$ happens to be in $M$, then $G(S(q_0),q_0)$ will be true in $M$. Now, we already know that each 
$S(q_0)\in M$. Then consider the following function which is definable over $M$:
\[ x\in\mathcal{G}\wedge G(y,x) \]
Because the domain, $\mathcal{G}$, is in $M$, so is the range, and because we have shown each $S(q_0)$ to be in $M$, the 
range is the set of all $S(q_0)$ (for $q_0\in\mathcal{G}$). Then the union of this, which is also in $M$, 
is just $S$ itself.

\subsection*{Theorem 2.3}
It is not fully clear to me why 
$p\vDash\lbc(\textrm{Ordinal}(f)\wedge\textrm{Ordinal}(\hat{\delta}))\supset (f\in\hat{\delta}\vee\hat{\delta}\in f\vee f\approx\hat{\delta})\rbc$.

A familiar error (see my notes on the previous chapter for what this error is) occurs in this proof.

\section*{Section 3}

\subsection*{Theorem 3.2}
A familiar error (see my notes on the previous chapter for what this error is) occurs in this proof.

\end{document}
